% =====================================================================
% Table 3c --- the closure matrix, transposed
% ---------------------------------------------------------------------
% Revised 2026-10-18: financing conditions now name the balancing variable
% (tau, B_gov, F); F is the endogenous net external transfer (ADR-0019) and
% the programme is booked as B_gov (F3 only). Household expenditure and the
% wage-income tax are defined once, in the caption.
% LAYOUT (2026-10-17, Table 2 convention):
%   * every block carries its signifier on its OWN line, with the condition
%     indented beneath it: \lblk (labour block), \fblk (financing block),
%     \pblk (the production pin, used inside the Leontief labour block).
%   * \lblk reserves a FIXED height \tbh, so the labour block occupies the
%     same vertical space in every cell and the "financing:" signifiers of
%     a row sit on the SAME baseline in all five columns (verified by
%     measuring the word boxes in the compiled PDF). The Leontief cell's
%     extra "production:" block lives inside that reserved height.
%   * \vspace{6pt} between the blocks = the padding before "financing:".
%   * LINE BREAKS INSIDE CELLS: cells are p{} columns, where a literal \\
%     is the ROW separator, NOT a line break. \newline is for the header
%     labels only: \lblk already ends its paragraph, so a \newline after it
%     dies with "There's no line here to end".
% \tbh is the single tunable: raise it if a condition grows a line. An
% "Overfull \vbox" warning in the log means it is too small.
% Block labels: labour market: the labour closure condition (the column);
% production: the substitution pin (only the Leontief column); financing:
% the financing condition (the row). British spelling ("labour") kept.
% =====================================================================
\begin{table}[htbp]
\centering
\footnotesize
\setlength{\tabcolsep}{2pt}
\setlength{\parindent}{0pt}
\setlength{\parskip}{0pt}
\renewcommand{\arraystretch}{1.3}
\def\tbh{38pt}% reserved height of the labour block: the tallest of the five
%   natural heights (the standard-CGE cell, 36.6pt) plus a half-point of slack.
%   Measured, not guessed: see the header comment.
\newcommand{\tagline}[1]{\textbf{#1:}}
\newcommand{\pblk}[2]{\tagline{#1}\par\hspace{0.6em}#2\par}
\newcommand{\lblk}[3]{\parbox[t][\tbh][t]{\linewidth}{\tagline{#1}\par\hspace{0.6em}#2\par#3}\par\vspace{6pt}}
\newcommand{\fblk}[2]{\tagline{#1}\par\hspace{0.6em}#2}
\begin{tabular}{@{}>{\raggedright\arraybackslash}p{1.85cm}
                  >{\raggedright\arraybackslash}p{2.45cm}
                  >{\raggedright\arraybackslash}p{2.45cm}
                  >{\raggedright\arraybackslash}p{2.45cm}
                  >{\raggedright\arraybackslash}p{2.45cm}
                  >{\raggedright\arraybackslash}p{2.45cm}@{}}
\toprule
\textbf{Financing closure}
& \textbf{No-reallocation}\newline \scriptsize immobile labour, sectoral wages ($\eta = 0$)
& \textbf{Standard CGE}\newline \scriptsize one flexible wage, mobile labour ($\eta = 1$)
& \textbf{Elastic labour supply}\newline \scriptsize behavioural supply on the real wage
& \textbf{Fixed real wage}\newline \scriptsize real wage pegged, employment uncapped
& \textbf{Leontief closure}\newline \scriptsize real wage pegged, quantities absorb \\
\midrule
\textbf{F1}\newline \scriptsize repurposing of consumption
& \lblk{labour market}{$\eta = 0$: $L_i^{\mathrm{cm}}(p, y_i, w_i) = \bar L_i$, sectoral $w_i$}{}
  \fblk{financing}{$B_{\mathrm{gov}} = 0$, $\tau = 0$, $\tilde\beta_i \propto \beta_i d_i$}
& \lblk{labour market}{$L_i = L_i^{\mathrm{cm}}(p, y_i, w)$, $\sum_i L_i = \bar L$, one wage $w$}{}
  \fblk{financing}{$B_{\mathrm{gov}} = 0$, $\tau = 0$, $\tilde\beta_i \propto \beta_i d_i$}
& \lblk{labour market}{$L_i^{\mathrm{cm}}(p, y_i, w_i) = \bar L_i \allowbreak \bigl(\tfrac{w_i/\Pi(p)}{w_i^{0}/\Pi^{0}}\bigr)^{\eta_{s,i}}$}{}
  \fblk{financing}{$B_{\mathrm{gov}} = 0$, $\tau = 0$, $\tilde\beta_i \propto \beta_i d_i$}
& \lblk{labour market}{$w/\Pi(p) = \bar w = 1$, $L_i = L_i^{\mathrm{cm}}(p, y_i, w)$}{}
  \fblk{financing}{$B_{\mathrm{gov}} = 0$, $\tau = 0$, $\tilde\beta_i \propto \beta_i d_i$}
& \lblk{labour market}{$w/\Pi(p) = 1$, $L_i = L_i^{\mathrm{cm}}(p, y_i, w)$}{\pblk{production}{$(\theta, \epsilon, \sigma) \to 0^{+}$}}
  \fblk{financing}{$B_{\mathrm{gov}} = 0$, $\tau = 0$, $\tilde\beta_i \propto \beta_i d_i$} \\
\midrule
\textbf{F2}\newline \scriptsize refinancing via taxes
& \lblk{labour market}{$\eta = 0$: $L_i^{\mathrm{cm}}(p, y_i, w_i) = \bar L_i$, sectoral $w_i$}{}
  \fblk{financing}{$\tau\, w \sum_i L_i = \sum_i p_i g_i$, $B_{\mathrm{gov}} = 0$}
& \lblk{labour market}{$L_i = L_i^{\mathrm{cm}}(p, y_i, w)$, $\sum_i L_i = \bar L$, one wage $w$}{}
  \fblk{financing}{$\tau\, w \sum_i L_i = \sum_i p_i g_i$, $B_{\mathrm{gov}} = 0$}
& \lblk{labour market}{$L_i^{\mathrm{cm}}(p, y_i, w_i) = \bar L_i \allowbreak \bigl(\tfrac{w_i/\Pi(p)}{w_i^{0}/\Pi^{0}}\bigr)^{\eta_{s,i}}$}{}
  \fblk{financing}{$\tau\, w \sum_i L_i = \sum_i p_i g_i$, $B_{\mathrm{gov}} = 0$}
& \lblk{labour market}{$w/\Pi(p) = \bar w = 1$, $L_i = L_i^{\mathrm{cm}}(p, y_i, w)$}{}
  \fblk{financing}{$\tau\, w \sum_i L_i = \sum_i p_i g_i$, $B_{\mathrm{gov}} = 0$}
& \lblk{labour market}{$w/\Pi(p) = 1$, $L_i = L_i^{\mathrm{cm}}(p, y_i, w)$}{\pblk{production}{$(\theta, \epsilon, \sigma) \to 0^{+}$}}
  \fblk{financing}{$\tau\, w \sum_i L_i = \sum_i p_i g_i$, $B_{\mathrm{gov}} = 0$} \\
\midrule
\textbf{F3}\newline \scriptsize external finance
& \lblk{labour market}{$\eta = 0$: $L_i^{\mathrm{cm}}(p, y_i, w_i) = \bar L_i$, sectoral $w_i$}{}
  \fblk{financing}{$B_{\mathrm{gov}} = \sum_i p_i g_i$, $\tau = 0$}
& \lblk{labour market}{$L_i = L_i^{\mathrm{cm}}(p, y_i, w)$, $\sum_i L_i = \bar L$, one wage $w$}{}
  \fblk{financing}{$B_{\mathrm{gov}} = \sum_i p_i g_i$, $\tau = 0$}
& \lblk{labour market}{$L_i^{\mathrm{cm}}(p, y_i, w_i) = \bar L_i \allowbreak \bigl(\tfrac{w_i/\Pi(p)}{w_i^{0}/\Pi^{0}}\bigr)^{\eta_{s,i}}$}{}
  \fblk{financing}{$B_{\mathrm{gov}} = \sum_i p_i g_i$, $\tau = 0$}
& \lblk{labour market}{$w/\Pi(p) = \bar w = 1$, $L_i = L_i^{\mathrm{cm}}(p, y_i, w)$}{}
  \fblk{financing}{$B_{\mathrm{gov}} = \sum_i p_i g_i$, $\tau = 0$}
& \lblk{labour market}{$w/\Pi(p) = 1$, $L_i = L_i^{\mathrm{cm}}(p, y_i, w)$}{\pblk{production}{$(\theta, \epsilon, \sigma) \to 0^{+}$}}
  \fblk{financing}{$B_{\mathrm{gov}} = \sum_i p_i g_i$, $\tau = 0$} \\
\bottomrule
\end{tabular}
\caption{The expanded closure matrix: three financing closures
(rows) crossed with five labour-market closures (columns). Each cell
states the labour-market condition (the column), the production pin where
one applies (the Leontief column only), and the financing condition (the
row). Common to all cells: household expenditure is
$E^{h} = (1-\tau)\,w\sum_i L_i + F$, where $\tau$ is a proportional
wage-income tax (nonzero only under F2) and $F$ is the endogenous net
external transfer that closes the external account,
$S + T_{\mathrm{int}} + M - (I+X) = F + B_{\mathrm{gov}}$. The programme
cost $B_{\mathrm{gov}}$ is booked externally (F3 only); $\sum_i p_i g_i$ is
valued at each cell's equilibrium prices. Under F2 the tax $\tau$ balances
the programme at each cell's wage bill; with one household this equals a
lump-sum transfer of $\sum_i p_i g_i$ out of household income. Under F2 and F3
the programme enters final demand as additive real demand $g_i$; under F1 it
is a reallocation of the household preference weights, $\tilde\beta_i \propto
\beta_i d_i$, with no additive demand. $\theta$, $\epsilon$ and $\sigma$ are the
substitution elasticities whose limit $0^{+}$ gives fixed coefficients.
Notation as in Table~\ref{tab:closure-conditions}.}
\label{tab:closure-matrix-3x5}
\end{table}
